\section{Postfix to Prefix Conversion}
\label{sec:sq-postfix-to-prefix}

\problemstatement{Given a postfix expression, convert it to prefix
notation.}

\begin{patternbox}
Identical setup to Section~\ref{sec:sq-postfix-to-infix} -- the only
change is \emph{how} two popped sub-expressions are combined: instead of
$(\textit{left}\ \textit{op}\ \textit{right})$, prefix writes the operator
\textbf{first}, with no parentheses needed at all (prefix, like postfix,
is unambiguous without them).
\end{patternbox}

\subsection*{Understanding the problem carefully}

The stack-of-sub-expressions technique from
Section~\ref{sec:sq-postfix-to-infix} carries over completely unchanged;
only the string-building step at each operator differs.

\begin{ideabox}[Combine as Prefix Instead of Infix]
On an operator $\textit{op}$: pop $\textit{right}$, pop $\textit{left}$
(same order as before); push $\textit{op} + \textit{left} +
\textit{right}$ (operator first, then both sub-expressions
back-to-back, no parentheses or spaces needed since prefix length is
self-delimiting when operands are single characters, or separated by a
delimiter for multi-character operands).
\end{ideabox}

\subsection*{Algorithm}

\begin{enumerate}
  \item Initialize an empty stack of strings.
  \item For each token: operand $\to$ push as a string. Operator
        $\textit{op}$ $\to$ pop $\textit{right}$, pop $\textit{left}$;
        push $\textit{op} + \textit{left} + \textit{right}$.
  \item The stack's single remaining element is the prefix result.
\end{enumerate}

\subsection*{Dry run}

\dryrun{Take $\texttt{"ab+c*"}$.}

\begin{itemize}
  \item \texttt{a}: push. Stack $=[\texttt{a}]$.
  \item \texttt{b}: push. Stack $=[\texttt{a},\texttt{b}]$.
  \item \texttt{+}: pop right$=\texttt{b}$, left$=\texttt{a}$. Push
        \texttt{+ab}. Stack $=[\texttt{+ab}]$.
  \item \texttt{c}: push. Stack $=[\texttt{+ab},\texttt{c}]$.
  \item \texttt{*}: pop right$=\texttt{c}$, left$=\texttt{+ab}$. Push
        \texttt{*+abc}.
\end{itemize}

Result: \texttt{*+abc}, correctly representing $(a+b)*c$ in prefix.

\subsection*{C++ implementation}

\begin{lstlisting}
#include <bits/stdc++.h>
using namespace std;

string postfixToPrefix(string s) {
    stack<string> st;

    for (char c : s) {
        if (isalnum(c)) {
            st.push(string(1, c));
        } else {
            string right = st.top(); st.pop();
            string left = st.top(); st.pop();
            st.push(string(1, c) + left + right);
        }
    }
    return st.top();
}

int main() {
    cout << postfixToPrefix("ab+c*") << endl;
    return 0;
}
\end{lstlisting}

\begin{complexitybox}
\textbf{Time Complexity.} $O(n^2)$ worst case with naive string
concatenation, as in Section~\ref{sec:sq-postfix-to-infix}.
\textbf{Space Complexity.} $O(n)$.
\end{complexitybox}

\begin{takeawaybox}
This section is intentionally a near-copy of
Section~\ref{sec:sq-postfix-to-infix} -- the lesson is that once you have
the ``stack of sub-expressions, combine on every operator'' skeleton
right, every output \emph{format} (parenthesized infix, prefix, or
postfix-of-a-different-source) is just a one-line change to how the
combination step builds its string.
\end{takeawaybox}
